\medskip

Soit $\left(u_n\right)$ la suite définie par $u_0 = 1$ et pour tout entier naturel $n$, \: 
$u_{n+1} = \dfrac{u_n}{1 + u_n}$.

\begin{enumerate}

\item 
	\begin{enumerate}
		\item $u_1=u_{n+1} = \dfrac{u_0}{1 + u_0}=\dfrac{1}{2}$, $u_2 = \dfrac{u_1}{1 + u_1}=\dfrac{1}{3}$  et $u_3 = \dfrac{u_2}{1 + u_2} = \dfrac{1}{4}$.


		\item On complète les lignes 3 et 6 du script python ci-dessous pour que liste(k) prenne en paramètre un entier naturel k et renvoie la liste des premières valeurs de la suite $\left(u_n\right)$ de $u_0$ à $u_k$.

\begin{center}
\begin{tabularx}{0.5\linewidth}{|l|X|}\hline
1.&def liste(k) :\\
2.&\qquad L = []\\
3.&\qquad u = {\blue 1}\\
4.&\qquad for i in range(0, k+1) :\\
5.&\quad \qquad L.append(u)\\
6.&\quad \qquad u = {\blue u/(1+u)}\\
7.&\qquad return(L)\\ \hline
\end{tabularx}
\end{center}
	\end{enumerate}
\item On admet que, pour tout entier naturel $n$, $u_n$ est strictement positif.

$u_{n+1} - u_n= \dfrac{u_n}{1 + u_n}-u_n=\dfrac{u_n-u_n(1+u_n)}{1+u_n}=\dfrac{-u_n^2}{1+u_n}<0$ car $u_n$ est strictement positif.

La suite $(u_n)$ est donc décroissante.

\item 
$\left(u_n\right) $ est décroissante et minorée par $0$; d'après le théorème de la convergence monotone, la suite  $(u_n)$ est convergente vers une limite $\ell$ positive ou nulle.

\item 
On considère la fonction $f(x)=\dfrac{x}{1+x}$ , on a alors $u_{n+1} = f\left(u_n\right)$.

La fonction $f$ est continue, car dérivable sur $\R^+$, donc la limite $\ell$ vérifie  l'égalité $f(\ell) = \ell$.

On résout l'équation: $\dfrac{\ell}{1+\ell}= \ell \iff \ell = \ell(1 + \ell) \iff 0= \ell ^2 \iff \ell = 0$.

Donc la suite $(u_n)$ a pour limite 0.

\item 
	\begin{enumerate}
		\item 
Au vu des premiers termes, on peut conjecturer que pour $n \in \N$, on a $u_n = \dfrac{1}{n+1}$
				
		\item On démontre par récurrence la conjecture précédente.
Soit $P_n$ la propriété: $u_n=\dfrac{1}{n+1}$.
		
$\bullet~~$Initialisation :
		
pour $n=0$, $u_0=\dfrac{1}{0+1}=1$ donc $P_0$ est vraie.

$\bullet~~$Hérédité : 
		
Supposons que $P_n$ est vraie pour $n \in \N$, on a donc $u_n=\dfrac{1}{n+1}$.
		
On a alors $u_{n+1}=\dfrac{u_n}{1+u_n}=\dfrac{\dfrac{1}{n+1}}{1+\dfrac{1}{n+1}}=\dfrac{\dfrac{1}{n+1}}{\dfrac{n+1+1}{n+1}}=\dfrac{1}{1+(n+1)}$ : la relation est vraie au rang $(n + 1)$.

$P_{n+1}$ est donc vraie.
		
Comme $P_0$ est vraie, et que pour $n \in \N$,\: $P_n$ vraie entraîne $P_{n+1}$ vraie, d'après le principe de récurrence  on peut dire que pour tout $n \in \N$, $P_n$ est vraie.
	\end{enumerate}
\end{enumerate}

\vspace{1cm}
